% !TEX root = ../manuscript.tex

\section{L'apprentissage automatique appliqué à la santé}

Ces dernières années, l'apprentissage automatique s'est imposé comme un outil incontournable dans le domaine de la santé. Des travaux pionniers comme ceux d'Obermeyer \& Emanuel~\cite{obermeyer2016} ont montré que les algorithmes de ML peuvent surpasser les méthodes statistiques classiques pour prédire des événements cliniques. Plus récemment, une revue publiée dans \textit{Nature Medicine}~\cite{rajpurkar2022} dresse un panorama des applications actuelles~: diagnostic médical assisté par IA, prédiction de risque de maladies chroniques, personnalisation des traitements.

Dans ce contexte, les jeux de données combinant indicateurs biologiques et habitudes de vie constituent un terrain d'étude particulièrement riche. L'idée centrale est que les maladies chroniques ne sont pas seulement le résultat de facteurs génétiques, mais résultent d'une accumulation de comportements quotidiens mesurables~\cite{booth2012}.

\section{Régression~: prédire des indicateurs biologiques}

La régression linéaire est l'une des méthodes les plus anciennes et les plus utilisées en statistique appliquée à la santé~\cite{james2021}. Pour pallier ses hypothèses fortes, des méthodes de \textbf{régression régularisée} ont été développées. La régression Ridge~\cite{hoerl1970} pénalise les coefficients trop grands via un terme $\lambda \|\mathbf{w}\|_2^2$, réduisant la variance du modèle. La régression Lasso~\cite{tibshirani1996} utilise une pénalisation $L_1$ qui force certains coefficients à zéro, réalisant ainsi une \textbf{sélection automatique de variables}.

Les méthodes ensemblistes comme les \textbf{forêts aléatoires}~\cite{breiman2001} et le \textbf{gradient boosting} (XGBoost~\cite{chen2016}) offrent des performances souvent supérieures sur des données tabulaires, en capturant les interactions non-linéaires entre variables.

\section{Classification~: prédire le risque de maladie}

La prédiction du risque de maladie est l'un des cas d'usage les plus étudiés en santé numérique. Des modèles de classification supervisée ont été appliqués avec succès à la prédiction du diabète de type 2~\cite{kavakiotis2017}, des maladies cardiovasculaires~\cite{mohan2019} ou encore de l'hypertension~\cite{olawade2023}.

Parmi les algorithmes les plus utilisés~:
\begin{itemize}
  \item La \textbf{régression logistique}~\cite{james2021}, modèle linéaire interprétable de référence.
  \item Les \textbf{forêts aléatoires}~\cite{breiman2001} et le \textbf{gradient boosting}~\cite{chen2016}, qui dominent les benchmarks sur données tabulaires.
\end{itemize}

Un point de vigilance important est le \textbf{déséquilibre des classes}~\cite{japkowicz2002}~: dans notre dataset, 75,2\% des individus sont sans risque contre seulement 24,8\% à risque. Un modèle naïf qui prédit toujours "sans risque" atteint déjà 75\% d'accuracy --- sans rien apprendre du tout.

Pour corriger ce biais, deux stratégies sont principalement utilisées. L'\textbf{ajustement des poids de classe} (\texttt{class\_weight='balanced'}) pénalise davantage les erreurs sur la classe minoritaire pendant l'entraînement, sans modifier les données. Le \textbf{SMOTE}~\cite{chawla2002} (\textit{Synthetic Minority Over-sampling Technique}) adopte une approche différente~: il crée synthétiquement de nouveaux exemples de la classe minoritaire en interpolant entre observations existantes. Pour un individu $x_i$ et l'un de ses $k$ plus proches voisins $x_{nn}$ dans la classe minoritaire~:

\begin{equation}
x_{\text{new}} = x_i + \lambda \cdot (x_{nn} - x_i), \quad \lambda \in [0, 1]
\end{equation}

Les deux stratégies sont appliquées et comparées dans ce projet. Comme on le verra, aucune ne parvient à améliorer significativement l'AUC~: cela confirme que le problème vient de l'absence de signal dans les données, et non du déséquilibre des classes.

\section{Clustering~: identifier des profils de mode de vie}

Le clustering vise à regrouper des individus similaires sans variable cible prédéfinie~\cite{macqueen1967}. Dans le domaine de la santé, cette approche est utilisée pour la \textbf{stratification de patients}~\cite{obermeyer2019}.

L'algorithme \textbf{K-Means}~\cite{macqueen1967} reste la méthode de référence. La \textbf{classification hiérarchique ascendante (CAH)} offre une représentation visuelle via le dendrogramme. L'évaluation d'un clustering se fait via le \textbf{coefficient de silhouette}~\cite{rousseeuw1987}~:

\begin{equation}
s(i) = \frac{b(i) - a(i)}{\max(a(i),\, b(i))}
\end{equation}

où $a(i)$ est la distance moyenne de l'observation $i$ aux autres points de son cluster, et $b(i)$ est la distance moyenne au cluster voisin le plus proche. Une valeur proche de 1 indique un clustering de qualité.

\section{Interprétabilité des modèles}

Un modèle performant ne suffit pas, notamment dans le domaine de la santé~: il faut aussi pouvoir \textbf{expliquer} ses prédictions~\cite{rudin2019}. Deux grandes familles d'approches se distinguent.

Les \textbf{modèles intrinsèquement interprétables} (régression linéaire, régression logistique, arbres de décision) permettent de lire directement les relations~: les coefficients d'une régression indiquent immédiatement le sens et l'amplitude de l'effet de chaque variable.

Les \textbf{méthodes post-hoc} sont appliquées \textit{après} l'entraînement, sur n'importe quel modèle --- y compris les boîtes noires comme Random Forest ou XGBoost. L'idée est simple~: ne pas sacrifier les performances prédictives pour avoir de l'interprétabilité, mais expliquer le modèle a posteriori. \textbf{SHAP}~\cite{lundberg2017} (\textit{SHapley Additive exPlanations}) est la méthode post-hoc la plus utilisée aujourd'hui. Issue de la théorie des jeux coopératifs, elle répond à la question~: \textit{quelle est la contribution équitable de chaque variable à cette prédiction précise ?} La valeur SHAP $\phi_j$ de la variable $j$ est~:

\begin{equation}
\phi_j = \sum_{S \subseteq F \setminus \{j\}} \frac{|S|!\,(|F|-|S|-1)!}{|F|!} \Bigl[ f(S \cup \{j\}) - f(S) \Bigr]
\end{equation}

où $F$ est l'ensemble des variables, $S$ un sous-ensemble ne contenant pas $j$, et $f(S)$ la prédiction du modèle avec uniquement les variables de $S$. Concrètement, pour chaque individu, SHAP produit un score signé par variable~: positif si la variable pousse vers le risque, négatif sinon.

Dans ce projet, SHAP est utilisé comme \textbf{complément visuel}, aux côtés des feature importances (Random Forest, XGBoost) et des coefficients de la régression logistique --- qui restent les interprétations principales.
